%*******************************************************
% Abstract (English) + Contents
%*******************************************************
\pdfbookmark[1]{Abstract}{abstract}

\chapter*{Abstract}

The digital transition is reshaping the relationship between economic systems and energy infrastructures. Yet the analytical frameworks used to assess this transformation remain fragmented: economic models tend to treat ICT as an exogenous productivity shock, while energy-climate models struggle to endogenize the growing material demands of computation. This thesis examines the economic and energetic dimensions of ICT infrastructure through a multidisciplinary approach combining history of innovation, structural economic analysis, and energy-climate modeling.

Chapter~1 traces the phenomenological history of ICT from the telegraph to the AI-powered datacenter, showing how economic and energy questions emerge naturally from the technical trajectory. Chapter~2 develops a systematic economic analysis of the hypothesis that computation constitutes a new structural regime. Chapter~3 confronts the promises of ICT-enabled decarbonization with empirical evidence on barriers, adoption failures, and systemic conditions. Chapters~4 through~6 implement this analysis within the IMACLIM modeling framework, constructing scenarios that endogenize the ICT-energy coupling. Chapter~7 discusses the policy implications.

The central argument is that ICT infrastructure functions as a \emph{pharmakon}---simultaneously entropy-producing and entropy-reducing---and that the datacenter, as the locus where GPT-Electricity and GPT-ICT converge, represents a structural break requiring new analytical tools.

\cleardoublepage

%*******************************************************
% Contents (English, detailed — style Pottier)
%*******************************************************
\pdfbookmark[1]{Contents}{contents}
\renewcommand{\contentsname}{Contents}
\setcounter{tocdepth}{2}
\tableofcontents